\section{Activity: Does Air Travel with Sound?}
\subsection{Materials required}
\begin{enumerate}
    \item A small piece of tissue paper
    \item A sound-producing object, such as a bell or a plate and spoon
\end{enumerate}
\subsection{Procedure}
Hold the piece of tissue paper about 20 to 30 cm away from the sound source. Produce a sound using the bell, or tap the plate gently with a spoon. Watch the tissue paper carefully. Repeat the activity with the tissue paper at different distances from the sound source. Observe the movement of the tissue paper.

\newthought{Observe:} Does the tissue paper move? Does it move as if air is flowing from the sound source towards it? How is this different from the movement of the tissue when you blow air towards it?

\newthought{Discuss:} When you blow air towards the tissue, air flows from your mouth to the tissue and makes it move. When you produce a sound, you do not see air flowing from the sound source towards the tissue in the same way. Instead, the air particles near the sound source move back and forth. They push and pull the particles next to them. The disturbance passes from one particle to the next and eventually reaches our ears. The air particles do not travel all the way from the sound source to our ears. They mainly move back and forth around their usual positions.

When you hear a sound, it is the disturbance, or vibration, that travels from the source to your ear. The air itself does not travel all the way from the source to your ear.